
This class represents an element, for example a flip-flop,
a gate or a block of RAM.


{\bf String cellname} - the element name to be used in the netlist, e.g.
"FDRS" or "AND2".

{\bf String ident} - a unique identifier string of the form "b123" (see
{\bf icount} below).


{\bf private static HashMap\verb'<'String, Ports\verb'>' celldecls} -  a
map of cell types. The map key is the cell name and the value is class {\bf
Ports} describing the input and output ports. Ports contains 3 sets,
inputs, outputs and 3-state outputs. Each of these sets contains Strings
giving the port names.


{\bf HashSet\verb'<'String\verb'>' inportset} - this is a set of the input
pins for the element. It will be a common set for the element type stored
in an instance of the {\bf Ports} class kept in a static map {\bf
celldecls} (see below).

{\bf HashSet\verb'<'String\verb'>' outportset} - like inportset above this
is a set of output pins for the element.

{\bf HashSet\verb'<'String\verb'>' tsportset} -  like inportset above this
is a set of 3-state output pins for the element.

The above sets are the same sets as those contained in the Ports class in
the 'celldecls' map indexed by the cell name of this instance. When a port
connection is added to an element the port name is checked against the
appropriate set and if not present is added. Since the port sets are shared
among elements with the same cell name the port sets for an element type
are built up by usage and are not pre-defined. This can be a problem where
nested gates are combined into a single gate during optimisation since that
element prototype might not have been encountered during netlist building.
For that reason AND, OR, XOR, XNOR gates and LUTs are initialised in map
{\bf celldecls} for number of inputs from 1 to the LUT width
(XTDECode.initialise())

{\bf HashMap\verb'<'String, Net\verb'>' inports} - a map
of input ports where the map key is the port name and the value is an
instance of class Net.

{\bf HashMap\verb'<'String, Net\verb'>' outports} -  a
map of 2-state output ports where the map key is the port name and the
value is an instance of class Net.

{\bf HashMap\verb'<'String, Net\verb'>' tsports} -  a map
of 3-state output ports where the map key is the port name and the value
is an instance of class Net.

Maps 'inports', 'outports' and 'tsports' represent the input and output
connections to the element. The keys to these maps are the port name
strings and the associated values are the connected nets (class 'Net').


{\bf Gtype element\_type} - a gate type enum used to identify elements
that can be optimised. The gate type enums are defined in class {\bf
NetConstants}.

{\bf private static ArrayList\verb'<'Element\verb'>' instances = new
ArrayList\verb'<'Element\verb'>'(20000)} - a list of all elements in the
netlist. Optimisable elements are added to the front of the list and a
count of these is kept in static variable {\bf ocount}.
Elements whose type is NONE are added at the end of the list, all others
being inserted at the front. Thus the elements that are labeled as being
subject to optimisation are located at the head of list 'instances'. The
number of such entries in the list is given by 'ocount'. This scheme avoids
traversing the whole list of elements when performing an optimisation pass.

{\bf String expression} - if a logical expression was used to generate a
property, for example an expression for a LUT, then that expression is
stored here. It is kept in this field for display purposes but is
translated into a hexadecimal string and stored in field 'property\_init'
for later inclusion in the netlist file as an INIT property to initialise
the element (LUT contents or flip-flop initial state).

{\bf String property\_init} - properties of the element. This might be
initialisation for a LUT or a flip-flop.

{\bf public static Element gnd} \\
{\bf public static Element vcc} -
publicly available GND and VCC elements. These are assigned in the
constructor XElements() where also publicly available Nets {\bf LO} and
{\bf HI} are created and connected to their output ports.

{\bf int count = 0} - this is a counter used for the schematic netlist
display to identify elements shown more than once in the schematic in
which case they are coloured to highlight the fact.

{\bf private static int ocount = 0} - count of optimisable elements -
see above.

{\bf private static int icount = 0} - unique integer for generating cell
identifiers (see {\bf ident} above). It is incremented after each use.

NOTE: when 3PL is run repeatedly via the GUI, class TDECode and all
subclasses down to the class specific to the FPGA family are
constructed afresh. Thus all static variables are re-created. Classes
Element and Net are NOT reconstructed and hence static variables in
these classes must be explicitly re-initialised on each run! In the case
of this class this is done by method init().


The following subsections describe constructors and methods. Many
simple methods which only provide read or write access to class
variables are omitted.

\subsection{constructors}

    \begin{lstlisting}[gobble=4, language=java]
        public Element (String name) { makeElement(name, null, Gtype.NONE); }
    \end{lstlisting}

    Construct a general element. These are given the gate type
    Gtype.NONE and are not subject to nelist optimisation. Parameter
    {\bf name} is the element name, e.g. "IBUF". The resulting element
    is given an arbitrary unique netlist block name of the form "b123".

    \begin{lstlisting}[gobble=4, language=java]
        public Element (String name, String bname) { makeElement(name, bname, Gtype.NONE); }
    \end{lstlisting}

    Construct a general element. These are given the gate type
    Gtype.NONE and are not subject to nelist optimisation. The block
    name must be unique. Parameter {\bf name} is the element name, e.g.
    "IBUF". Parameter {\bf bname} is the netlist block name.

    \begin{lstlisting}[gobble=4, language=java]
        public Element (String name, Gtype element_type) { makeElement(name, null, element_type); }
    \end{lstlisting}

    Construct a gate element. The gate type is used to direct netlist
    optimisation. Parameter {\bf name} is the element name, e.g. "IBUF".
    Parameter {\bf element\_type} is the element gate type, e.g.
    Gtype.MUXCY The resulting element is given an arbitrary unique
    netlist block name of the form "b123".

    Each of these constructors calls the common private method
    {\bf makeElement()}. Each element is added to the static list
    {\bf instances}, optimisable elements at the front of the list
    to avoid traversing the whole list during optimisation.

\subsection{method init()}

    \begin{lstlisting}[gobble=4, language=java]
        public static void init () {
            icount = 0; // unique integer for identifiers
            ocount = 0; // count of optimisable elements
            celldecls = new HashMap<String, Ports>();
            instances = new ArrayList<Element>(20000);
        }
    \end{lstlisting}

    This method initialises static variables. This must be done
    by a method, rather than during class initialisation, to
    allow code generation to be restarted when 3PL is repeatedly run
    interactively from the GUI.

\subsection{method putCelldecl()}

    There are two signatures for this.

    \begin{lstlisting}[gobble=4, language=java]
        public static void putCelldecl (String name, int nports) {
            celldecls.put(name, gatePorts(nports));
        }
    \end{lstlisting}

    This method puts a logic gate cell declaration in the cell map.
    Normally the cell map is filled during calls to the constructors
    for class {\bf Element}, but gate elements are entered explicitly
    to ensure that gate of all possible widths are present. This is
    because optimisation might reduce say an AND3 gate to an AND2 gate
    but the latter may not have been encountered in calls to the
    constructor.

    \begin{lstlisting}[gobble=4, language=java]
        public static void putCelldecl (String name, ArrayList inports, ArrayList outports) {
            .
            .
            .
        }
    \end{lstlisting}

    This method puts a general cell declaration in the cell map. This
    is used for elements specific to particular FPGAs, e.g. MULT18X18
    or DSP48E and ensures that all inputs and outputs are included in
    the definition even if they are not included during a call to
    a constructor.

\subsection{method gatePorts()}

    \begin{lstlisting}[gobble=4, language=java]
        public static Ports gatePorts (int n) {
            Ports   p = new Ports();
            p.outports.add("O");
            for (int i=0 ; i<n ; i++)
                p.inports.add("I" + i);
            return(p);
        }
    \end{lstlisting}

    This method creates a port class for a gate element. The parameter
    indicates the number of input ports and there is one output port.

\subsection{adding inputs}

    To add a single Net to a named port -
    \begin{lstlisting}[gobble=4, language=java]
        public void addInput (String port, Net n) {
            switch (element_type) {
            case AND:
            case OR:
            case XOR:
                break;
            default:
                inportset.add(port);
            }
            if (n == null)
                return;
            inports.put(port, n);
            n.connect(this, port, PortType.INPORT);
        }
    \end{lstlisting}
    AND, OR and XOR gates have varying numbers of inputs and are subject
    to changes to the number of inputs during optimisation, so the input
    port set should not be touched. It can be assumed that the input
    port sets for all these gates that are within the accepted number of
    inputs for the FPGA family will have been already created. For all
    others the port string is added to the set of inports for this
    element type. The last statement creates a reciprocal link from the
    Net to this element.

    To add an array of Nets -
    \begin{lstlisting}[gobble=4, language=java]
        public void addInputArray (String port, Net[] n) {
            if (n == null)
                return;
            for (int i=0 ; i<n.length ; i++)
                addInput(port + i, n[i]);
        }
    \end{lstlisting}
    The port string has "0", "1" etc. appended.

    To add an array of Nets with most significant bit padding -
    \begin{lstlisting}[gobble=4, language=java]
        public void addInputArray (String port, Net[] n, int size, Net pad) {
            int portnum = 0;
            int asize = (n == null) ? 0 : n.length;
            for (int i=0 ; i<asize ; i++)
                addInput(port + portnum++, n[i]);
            for (int i=asize ; i<size ; i++)
                addInput(port + portnum++, pad);
        }
    \end{lstlisting}
    The port string has "0", "1" etc. appended. If the array is smaller
    than the designated size the remaining inputs are connected to
    Net {\bf pad}. If the array is null all inputs are connected to
    Net {\bf pad}.

    To add an array of Nets with least significant bit padding -
    \begin{lstlisting}[gobble=4, language=java]
        public void addInputArrayTop (String port, Net[] n, int size, int max, Net pad) {
            int npad = max - size;
            int asize = (n == null) ? 0 : n.length;
            int portnum = 0;
            for (int i=0 ; i<npad ; i++)
                addInput(port + portnum++, pad);
            for (int i=0 ; i<size ; i++)
                addInput(port + portnum++, (i < asize) ? n[i] : Net.LO);
        }
    \end{lstlisting}
    If the Net array is smaller than {\bf size} then the most
    significant bits are padded as well.

    To add an array of Nets across two different ports -
    \begin{lstlisting}[gobble=4, language=java]
        public void addInputArray (
            String  port1,
            String  port2,
            Net[]   n,
            int     size1
        ) {
            int portnum1 = 0;
            int portnum2 = 0;
            int asize = (n == null) ? 0 : n.length;
            for (int i=0 ; i<asize ; i++)
                if (i < size1)
                    addInput(port1 + portnum1++, n[i]);
                else
                    addInput(port2 + portnum2++, n[i]);
        }
    \end{lstlisting}
    Parameter {\bf port1} is the first input port identifying string
    head. Parameter {\bf port2} is the second input port identifying
    string head. Parameter {\bf n} is the net array. Parameter {\bf
    size1} is the number of ports using the first head string.

    To add an array of Nets across two different ports with most
    significant bit padding -
    \begin{lstlisting}[gobble=4, language=java]
        public void addInputArray (
            String  port1,
            String  port2,
            Net[]   n,
            int     size1,
            int     size,
            Net     pad
        ) {
            int portnum1 = 0;
            int portnum2 = 0;
            int asize = (n == null) ? 0 : n.length;
            for (int i=0 ; i<asize ; i++)
                if (i < size1)
                    addInput(port1 + portnum1++, n[i]);
                else
                    addInput(port2 + portnum2++, n[i]);
            for (int i=asize ; i<size ; i++)
                if (i < size1)
                    addInput(port1 + portnum1++, pad);
                else
                    addInput(port2 + portnum2++, pad);
        }
    \end{lstlisting}



\subsection{adding outputs}

    For output Net and Net arrays these follow the same pattern as
    for inputs above -
    \begin{lstlisting}[gobble=4, language=java]
        public void addOutput (String port, Net n) {
            outportset.add(port);
            if (n == null)
                return;
            if (n.getNumOutputs() != 0)
                throw new ExEx("net 2nd source");
            outports.put(port, n);
            n.connect(this, port, PortType.OUTPORT);
        }
    \end{lstlisting}
    There is a builtin check here to ensure that the Net is not
    connected to any other element outputs. The last statement creates a
    reciprocal link from the Net to this element.

    And similarly -
    \begin{lstlisting}[gobble=4, language=java]
        public void addOutputArray (String port, Net[] n)
        public void addOutputArray (String port1, String port2, int size1, Net[] n)
    \end{lstlisting}

    For 3-state outputs -
    \begin{lstlisting}[gobble=4, language=java]
        public void addTS (String port, Net n) {
            tsportset.add(port);
            if (n == null)
                return;
            tsports.put(port, n);
            n.connect(this, port, PortType.TSPORT);
        }
    \end{lstlisting}


\subsection{method removeNet}

    This method removes a Net connected to an Element. The argument
    is the port to which the Net is connected.
    \begin{lstlisting}[gobble=4, language=java]
        public void removeNet (String port) {
            if (inports.containsKey(port)) {
                inports.get(port).disconnect(this, port);
                inports.remove(port);
            } else if (outports.containsKey(port)) {
                outports.get(port).disconnect(this, port);
                outports.remove(port);
            } else if (tsports.containsKey(port)) {
                tsports.get(port).disconnect(this, port);
                tsports.remove(port);
            } else
                throw new ExEx("net not connected to element");
        }
    \end{lstlisting}
    The code checks the input, output and 3-state port maps in turn
    looking for the port name string as key. Inputs are searched first
    as they are the more numerous. The get() method returns the Net
    and the disconnect() method removes the Element from the Net.
    The remove() then removes the key and associated value ()Net from
    the particular Element port map.


\subsection{method remapInputs}

    This method re-names the input ports of a gate Element. It is
    required when one or more inputs have been removed as a result
    of optimisation.
    \begin{lstlisting}[gobble=4, language=java]
        public void remapInputs () {
            int         n = inports.size();
            HashMap     ports_new = new HashMap();
            Net         inet;
            int         i = 0;
            for (Map.Entry me : inports.entrySet()) {
                String      newport = "I" + i++;
                String      port = (String)me.getKey();
                inet = (Net)me.getValue();
                inet.changePort(this, port, newport);
                ports_new.put(newport, inet);
            }

            inports = ports_new;
            cellname = element_type.toString() + n;

            Ports p;
            if (celldecls.containsKey(cellname)) {
                p = celldecls.get(cellname);
                inportset = p.inports;
                outportset = p.outports;
                tsportset = p.tsports;
            } else {
                // Remapped wide gates may not have declared cell
                // set entries - dont worry as wide gates will
                // be converted so will never need these sets anyway.
                inportset = null;
                outportset = null;
                tsportset = null;
            }
        }
    \end{lstlisting}
    The input port map is traversed and the Nets are placed in a new map
    where the keys (port names) are allocated sequentially as "I0",
    "I1" etc. That new map is then assigned as a replacement input port
    map. The Element cell name is changed to reflect the ne winput count,
    e.g. "AND5" may now be "AND3". The new cell name is now checked
    in the cell declaration map and if present the port sets for the
    Element are assigned from that map. If not present it means that
    the gate is a wide gate which has not been entered in
    the cell declaration map since it will be converted to use available
    library elements later.

\subsection{method changeElement()}

    This method changes an element. This is used during optimisation to
    change an element into another. Previous connections are optionally
    stripped. Parameter {\bf name} is the new element name, e.g. "INV".
    Parameter {\bf element\_type} is the element gate type, e.g.
    Gtype.INV. Parameter {\bf strip} if true results in the element being
    stripped of input and output Nets. An example might be an XOR2 gate
    with one input connected to VCC being changed to an INV of the
    other input.

    \begin{lstlisting}[gobble=4, language=java]
        public void changeElement (String name, Gtype element_type, boolean strip) {
            Ports   p;
            if (strip)
                strip();
            cellname = name;
            this.element_type = element_type;
            if (celldecls.containsKey(name))
                p = celldecls.get(name);
            else {
                p = new Ports();
                celldecls.put(name, p);
            }
            if (strip) {
                inports = new HashMap<String, Net>();
                outports = new HashMap<String, Net>();
                tsports = new HashMap<String, Net>();
            }
            inportset = p.inports;
            outportset = p.outports;
            tsportset = p.tsports;
        }
    \end{lstlisting}
    The Element is first optionally stripped of all inputs and outputs,
    depending on boolean argument {\bf strip}. The cell name string
    and element type are then overwritten from the supplied arguments.
    If the new element is contained in the cell declaration map then
    the ports sets are obtained from that map, otherwise new sets are
    created for the named cell. If the element has been stripped of
    connections then new port maps are created, otherwise the existing
    port maps remain. The port sets are then assigned from the sets
    obtained from the cell declarations map.

\subsection{method strip()}

    This method disconnects the Element from all input and output
    nets and then clears all input and output port maps. This is used
    prior to deleting an Element or changing all its port
    connections during optimisation.
    \begin{lstlisting}[gobble=4, language=java]
        public void strip () {
            for (Map.Entry me : inports.entrySet()) {
                String      pinname  = (String)me.getKey();
                Net         net = (Net)me.getValue();
                net.disconnect(this, pinname);
            }
            inports.clear();
            for (Map.Entry me : outports.entrySet()) {
                String      pinname  = (String)me.getKey();
                Net         net = (Net)me.getValue();
                net.disconnect(this, pinname);
            }
            outports.clear();
            for (Map.Entry me : tsports.entrySet()) {
                String      pinname  = (String)me.getKey();
                Net         net = (Net)me.getValue();
                net.disconnect(this, pinname);
            }
            tsports.clear();
        }
    \end{lstlisting}

\subsection{method addExpression()}

    This method adds an expression to a LUT. The expression string
    is kept for display purposes and is also translated into a
    property value for LUT initialisation. Parameter {\bf s} is the
    expression string.
    \begin{lstlisting}[gobble=4, language=java]
        public void addExpression (String s) {
            expression = s;
            property_init = Functions.bit_pattern(s, TDECode.lutwidth);
        }
    \end{lstlisting}

\subsection{method addProperty()}

    This method adds a property for LUT initialisation. Parameter
    {\bf s} is the property string.
    \begin{lstlisting}[gobble=4, language=java]
        public void addProperty (String s) {
            if (s == null)
                return;
            if (property_init != null)
                throw new ExEx("element INIT property multiply defined");
            property_init = s;
        }
    \end{lstlisting}



\subsection{methods for writing to the netlist file}

    Method
    \begin{lstlisting}[gobble=4, language=java]
        static public void outputCells (PrintWriter pw, String view)
    \end{lstlisting}
    outputs the cell port definitions to the EDIF netlist.
    Only cells which have been instantiated are written out.
    Some cells may have been optimised out such that no instances
    remain in which case they are not written to the EDIF file cell
    definition section. Parameter {\bf pw} is the output interface to
    the EDIF netlist file and Parameter {\bf view} is the EDIF view
    parameter (the string "view\_1" seems to be required and this
    originates in codegen.TDEList.ncode()).

    Method
    \begin{lstlisting}[gobble=4, language=java]
        static public void outputInstances (PrintWriter pw, String view)
    \end{lstlisting}
    outputs all Elements to the EDIF netlist.

    Method
    \begin{lstlisting}[gobble=4, language=java]
        static public void outputConstraints (String group, PrintWriter pw)
    \end{lstlisting}
    outputs all element constraints for a group. Parameter {\bf pw} is
    the output interface to the constraint file for the group (name.ncf
    for group "n" and name.ucf for group "u"). Note that this method is
    called from netlist.xilinx.XTDECode() which also writes net
    constraints and general constraints to constraint files.
